\documentclass[11pt,a4paper]{article}
\usepackage{amsmath,amssymb,amsthm}
\usepackage{graphicx}
\usepackage[colorlinks]{hyperref}
\newtheorem{theorem}{Theorem}
\newtheorem{definition}{Definition}
\newcommand{\normvec}[1]{\lVert #1 \rVert}

\title{A Test Document for Texpile}
\author{Texpile QA \and A Second Author}
\date{\today}

\begin{document}

\maketitle

\begin{abstract}
	This document exists to be compiled, folded, searched, spell-checked, and jumped around
	in. It covers sectioning, cross-references, citations, floats, tables, theorem
	environments, and listings. If any of it renders wrongly, that is the bug.
\end{abstract}

\section{Sectioning and cross-references}
\label{sec:intro}

This is the first section, and the next one is about mathematics. Sectioning commands should
colour as commands, and the titles should stay plain text. Here is a footnote.\footnote{The
footnote text should be spell-checked like any other prose.} Here is a
\href{https://texpile.com}{link with text}, and a bare one: \url{https://typst.app}.

\subsection{A subsection}
\label{sec:sub}

Deep nesting exists so the outline and the fold column have something to work with.

\paragraph{A paragraph heading} runs into its text, which is a different shape again.

\section{Mathematics}
\label{sec:math}

Inline math like $E = mc^2$ and $\alpha \neq \beta$ sits in the middle of a sentence, as does
the alternative form \( \sum_{k=1}^{n} k = \tfrac{n(n+1)}{2} \). The math preview should render
either on hover.

Display math, unnumbered:
\[
	\int_{0}^{\infty} e^{-x^2} \, dx = \frac{\sqrt{\pi}}{2}
\]

A numbered equation, which the text refers back to:
\begin{equation}
	\label{eq:euler}
	e^{i\pi} + 1 = 0
\end{equation}

Aligned equations, each numbered, using the custom macros from the preamble:

\begin{align}
\normvec{x}^2 & =\sum_{i=1}^{n}x_{i}^2 \label{eq:norm} \\
\operatorname*{\mathrm{arg\,min}}_{\theta\in\mathbb{R}^{d}}f(\theta) & \in\mathbb{R}^{d} \label{eq:argmin}
\end{align}

Both equations are reachable by reference; the middle line is deliberately unnumbered.

\begin{theorem}[Pythagoras]
	\label{thm:pythagoras}
	In a right triangle, $a^2 + b^2 = c^2$.
\end{theorem}

\begin{proof}
	Left to the reader, as is traditional.
\end{proof}

\begin{definition}
	\label{def:cauchy}
	A sequence $(x_n)$ is Cauchy if for every $\epsilon > 0$ there is $N \in \mathbb{N}$ such that
	$|x_n - x_m| < \epsilon$ for all $n, m > N$.
\end{definition}

\section{Lists}
\label{sec:lists}

\begin{itemize}
	\item A bullet.
	\item Another, with \emph{emphasis}, \textbf{bold}, \texttt{monospace}, and
	      \underline{underline}.
	\begin{itemize}
		\item A nested bullet.
	\end{itemize}
\end{itemize}

\begin{enumerate}
	\item First.
	\item Second.
	\begin{enumerate}
		\item Nested and renumbered.
	\end{enumerate}
\end{enumerate}

\begin{description}
	\item[Term] its definition.
	\item[Another term] a second definition.
\end{description}

\section{Verbatim}
\label{sec:code}

Inline verbatim: \verb|\newcommand{\foo}{bar}|. A verbatim block, where nothing should be
interpreted or spell-checked:

\begin{verbatim}
for i in range(10):
    print(f"{i}: {i ** 2}")   % this % is not a comment
\end{verbatim}

\section{Citations}
\label{sec:cites}

A single citation \cite{knuth1984}, a page-qualified one \cite[p.~42]{lamport1994}, and
several at once \cite{mittelbach2004, typst2023, curry1930}. A citation to an online source
sits here too \cite{unicode2024}. The picker should offer all of them, and each should
resolve to the author and year rather than the raw key.

\section{Things that commonly break}
\label{sec:edge}

Special characters: 100\% of \$5, a\_b, \{braces\}, \#hash, \&ampersand, and a tilde\textasciitilde{}.

Accents and non-ASCII, which the spell checker should not choke on: naïve, café, Gödel,
Erdős, Straße, ¿por qué?, 日本語, العربية, עברית.

An em-dash --- like this --- and an en-dash 1--10, plus ``smart quotes'' and `single' ones.

A line break forced with a double backslash: \\
this text is on the next line.

\bibliographystyle{plain}
\bibliography{refs}

\end{document}
